\begin{abstract}
Unsafe tool calls can accompany textual refusals. PRISM tests activation intervention transfer across precision, languages, and scales, separating construction, layer and dose selection, and evaluation by scenario family. BF16 activations weighted by BF16 or NF4 labels yield directions increasing unsafe calls by 76.5–78.2 percentage points (pp) in non-thinking Qwen3.5-9B, 6.3–6.9 in Ministral-3-8B, and 17.9–19.7 in Llama-3.1-8B; negative doses reduce rates by 5.6–6.0 and 25.1–26.5 pp in Ministral and Llama (Qwen lacked eligible state-changing calls for suppression). Maps learned from paired benign activations transfer directions to larger models, increasing unsafe calls by 76.01, 3.62, and 27.50 pp in Qwen3.5-27B, Ministral-3-14B, and Llama-3.1-70B, and reducing them by 7.65 and 35.15 pp in the latter two, with at most 2.84 pp loss on tool-required BFCL. In held-out Arabic, five-language NF4-supervised directions reduce Llama-3.1-8B violations by 21.95 pp; Ministral effects remain near zero. NF4 more often enables estimation with few scenarios but shows no clear intervention advantage. Llama-3.1-70B suppression spans 12.38–38.62 pp across three evaluators, and fixed interventions on executable SODA tasks also change unsafe-trajectory rates. Results demonstrate transfer across three families and its limits.
\end{abstract}
